\documentclass[14pt,a4paper]{extarticle}
\usepackage[T2A]{fontenc}
\usepackage[utf8]{inputenc}
\usepackage[english,russian]{babel}
\usepackage{lmodern}
\usepackage{geometry}
\usepackage{array}
\usepackage{tabularx}
\usepackage{longtable}
\usepackage{ltablex}
\usepackage{booktabs}
\usepackage{hyperref}
\usepackage{indentfirst}
\usepackage{microtype}
\usepackage{graphicx}
\usepackage{setspace}
\usepackage{listings}
\usepackage{graphicx}
\usepackage{caption}
\usepackage{float}
\usepackage{tocloft}
\usepackage{titlesec}
\usepackage{newtxtext,newtxmath}

\titleformat{\section}        % форматируем \section
  {\normalfont\bfseries}      % стиль: обычный + жирный
  {\thesection.}              % номер секции с точкой
  {2mm}                       % расстояние между номером и текстом
  {}                           % перед заголовком можно добавить что-то
\titlespacing*{\section}{1.5cm}{*3}{*1}  % {отступ слева}{отступ сверху}{отступ снизу}

\titleformat{\subsection}        % форматируем \section
  {\normalfont\bfseries}      % стиль: обычный + жирный
  {\thesubsection.}              % номер секции с точкой
  {2mm}                       % расстояние между номером и текстом
  {}                           % перед заголовком можно добавить что-то
\titlespacing*{\subsection}{1.5cm}{*3}{*1}  % {отступ слева}{отступ сверху}{отступ снизу}

\captionsetup{justification=centering}
\renewcommand{\ttdefault}{cmtt}
\keepXColumns
\setlength{\parindent}{3em}
\newgeometry{left=2.5cm,right=2.5cm,top=3cm,bottom=2.5cm}

\renewcommand{\contentsname}{\hspace*{4em}{\smallСодержание}}
\setlength{\cftbeforesecskip}{0.5em}
\renewcommand{\cftsecleader}{\cftdotfill{\cftdotsep}}
\renewcommand{\cftsecaftersnum}{.\hspace{2mm}}

\title{Отчёт по "Управление ИТ-проектами"}
\author{Выполнил студент группы М-РИТ-25-1, Станиславчук Сергей}

\begin{document}

\setstretch{1.5}
\noindent
\begin{minipage}[t]{0.01\textwidth}
    \includegraphics{~/programmer/++Programmer/5_1/LSTU.jpg}
\end{minipage}%
\begin{minipage}[c]{1.1\linewidth}
    \centering
    {\small
    МИНИСТЕРСТВО НАУКИ И ВЫСШЕГО ОБРАЗОВАНИЯ РФ \\[1em]
    ФЕДЕРАЛЬНОЕ ГОСУДАРСТВЕННОЕ БЮДЖЕТНОЕ \\
    ОБРАЗОВАТЕЛЬНОЕ УЧРЕЖДЕНИЕ ВЫСШЕГО ОБРАЗОВАНИЯ \\
    «ЛИПЕЦКИЙ ГОСУДАРСТВЕННЫЙ ТЕХНИЧЕСКИЙ УНИВЕРСИТЕТ»
    }
\end{minipage}%

\vspace{3em}

\noindent
\hspace{1em}%
\begin{minipage}[t]{0.8\textwidth}
\textbf{Институт} \hspace{2em} компьютерных наук \\
\textbf{Кафедра} \hspace{2em} автоматизированных систем управления
\end{minipage}


\vspace{6em}

\begin{center}
\textbf{Практическая работа №9} \\
по дисциплине \\ <<Управление ИТ-проектами>>
\end{center}

\vspace{6em}

\noindent
\begin{tabular*}{\textwidth}{@{\extracolsep{\fill}} l c r }
    Студент М-РИТ-25-1 & \makebox[4cm]{\hrulefill} & Станиславчук С.М. \\[2em]
Руководитель & \\
К.т.н., доцент & \makebox[4cm]{\hrulefill} & Ведищев В.В. \\
\end{tabular*}

\vfill
\begin{center}
    Липецк, 2026
\end{center}

\thispagestyle{empty}

\newpage
\tableofcontents
\newpage

\section{Анализ рисков проекта}

\subsection{Понятие риска}

Риск - функция, включающая в себя вероятность неблагоприятного события и его последствий.

Если риск равен 3,5\%, это означает, что вероятность неблагоприятного исхода составляет 3,5\% (или 0,035 в долях). Это сравнительно низкий риск, но его все же стоит учитывать при планировании, особенно если потери в случае неудачи будут значительными.

Риск реализации проекта — это вероятность возникновения проблем или неудач, которые могут повлиять на успешное завершение проекта. Эти риски могут быть связаны с техническими, финансовыми, организационными или внешними факторами.

\subsection{Характеристики риска}

Каждый риск характеризуется рядом параметров, позволяющих оценить его значимость. К основным характеристикам относятся вероятность возникновения, степень воздействия на проект, а также обнаруживаемость. В контексте данного проекта важными являются также такие характеристики, как влияние на безопасность данных, производительность системы и доступность сервиса внутри корпоративной сети.

\subsection{Наиболее распространённые риски в области IT}

Для разрабатываемой системы можно выделить ряд типовых рисков, характерных для IT-проектов. К ним относятся ошибки проектирования, уязвимости безопасности, сбои сетевого взаимодействия, проблемы масштабирования, а также внешние факторы, включая изменения законодательства.

\subsection{Оценка рисков}

Оценка рисков проводится на основе вероятности их возникновения и степени влияния на систему.

\subsection{Качественный анализ рисков}

Качественный анализ направлен на экспертную оценку рисков с точки зрения их влияния на систему и вероятности возникновения. Данный подход позволяет определить приоритеты обработки рисков и выделить наиболее критичные из них, требующие первоочередного внимания.

В рамках проекта используются методы качественного анализа, основанные на экспертных оценках и структурированном анализе факторов. Метод экспертных оценок предполагает привлечение участников команды разработки, обладающих практическим опытом в соответствующих областях. Дополнительно применяется SWOT-анализ, позволяющий рассмотреть риски в контексте внутренних и внешних факторов.

В качественном анализе рисков фиксируются ключевые параметры: возможное событие, момент его возникновения, последствия и вероятность реализации.

\begin{table}[h!]
\centering
\begin{tabular}{|p{4cm}|p{3cm}|p{5cm}|p{2cm}|}
\hline
\textbf{Событие} & \textbf{Когда возникает} & \textbf{Последствия} & \textbf{Вероятность} \\
\hline
Уязвимости безопасности & На этапе реализации и эксплуатации & Утечка данных, компрометация системы & Высокая \\
\hline
Сбои сетевого взаимодействия & При передаче данных & Потеря сообщений, задержки голосовой связи & Средняя \\
\hline
Ошибки архитектуры & На этапе проектирования & Сложность масштабирования, переработка системы & Средняя \\
\hline
Сложность внедрения & На этапе развёртывания & Увеличение сроков запуска системы & Высокая \\
\hline
Низкая производительность & На этапе эксплуатации & Снижение качества работы системы & Средняя \\
\hline
Изменения законодательства & В процессе эксплуатации & Необходимость переработки системы & Низкая \\
\hline
\end{tabular}
\caption{Параметры качественного анализа рисков}
\end{table}

Таким образом, качественный анализ позволяет структурировать информацию о рисках и определить, какие из них требуют первоочередного внимания при разработке системы.
Качественный анализ направлен на экспертную оценку рисков с точки зрения их влияния на систему и вероятности возникновения. Данный подход позволяет определить приоритеты обработки рисков и выделить наиболее критичные из них, требующие первоочередного внимания.

В рамках проекта используются методы качественного анализа, основанные на экспертных оценках и структурированном анализе факторов.

В предыдущей работе применялся SWOT-анализ, позволяющий рассмотреть риски в контексте сильных и слабых сторон проекта, а также внешних факторов. В ходе ранее проведённого SWOT-анализа были выявлены ключевые риски, связанные с изменением государственных законов, уязвимостями безопасности, сложностью внедрения и конкуренцией со стороны существующих решений.

Например, сильной стороной проекта является децентрализованная архитектура, обеспечивающая независимость от внешних сервисов. Однако это же приводит к увеличению сложности развёртывания и администрирования системы, что формирует соответствующий риск. Возможностью является рост спроса на защищённые корпоративные решения, в то время как угрозой выступает наличие зрелых конкурирующих платформ и возможные изменения в законодательстве.

\subsection{Количественный анализ рисков}

Количественный анализ рисков направлен на оценку потенциальных потерь проекта в случае реализации неблагоприятных событий. В отличие от качественного анализа, он позволяет выразить последствия рисков в измеримых величинах, таких как время задержки разработки, снижение производительности системы или затраты на устранение последствий.

В рамках разработки системы корпоративного обмена сообщениями количественный анализ применяется к рискам, имеющим высокую вероятность возникновения и значительное влияние на функционирование системы. Например, реализация уязвимости безопасности может привести к компрометации корпоративных данных. В таком случае потребуется проведение аудита, исправление уязвимостей и возможная переработка механизмов защиты, что приведёт к значительным временным затратам и увеличению сроков разработки.

Другим примером является риск сбоев сетевого взаимодействия. Потеря пакетов или высокая задержка при передаче голосовых данных может потребовать доработки сетевых протоколов и оптимизации передачи данных, что также увеличивает трудозатраты и сроки реализации проекта.

На основе результатов качественного и количественного анализа риски разделяются на три группы. К первой группе относятся риски с высокой вероятностью и значительным влиянием, такие как уязвимости безопасности и ошибки архитектуры. Для них требуется реализация мер защиты на ранних этапах разработки.

Ко второй группе относятся риски с высокой вероятностью, но меньшим влиянием, либо редкие, но критические события. К ним можно отнести проблемы производительности или изменения внешних условий. Такие риски необходимо контролировать и учитывать при планировании.

К третьей группе относятся риски с низкой вероятностью и незначительным влиянием, например отдельные ошибки пользовательского интерфейса. Они учитываются, но не требуют значительных ресурсов на предотвращение.

Таким образом, количественный анализ позволяет определить приоритеты управления рисками и рационально распределить ресурсы проекта, сосредоточив внимание на наиболее критичных угрозах.

\end{document}
